\section{데이터 형식과 세부 실행 규칙}
\label{app:protocol-details}

이 부록은 \ref{sec:system}–\ref{sec:audit}절의 운영 가정과
사용하는 값의 정확한 표현, 실행 조건을 정리한다.

\subsection{참여 주체와 감사키}
\label{app:system-model}
\label{app:audit-key}

등록 승인자는 최초 값과 동일 물량의 중복 등록 여부를 판단하고,
정책 관리자는 가공·발급 조건과 검증키, 기업의 가공 권한을 등록한다.
이 두 역할의 판단·등록과 감사자의 조사 시작점 선택·상태 집행을 신뢰한다.
감사자와 상태 관리자는 같은 주체이며, 감사 클라이언트는 그 주체의
소프트웨어로 별도의 신뢰 주체가 아니다.

원장은 확정 순서·변경되지 않는 수락 기록·인증된 상태 읽기와
원자적 거래를 제공한다고 가정한다. 원격 조회 결과도 인증해야 한다.
질량·재활용 크레딧·탄소와 운송·공정 증분은 외부에서 제공되며,
실제 측정·공정 수행과 제품·문서의 대응은 외부 확인에 의존한다.

송신자는 이전 증표의 개방값을 보관하고, 이전 거래 제출 전이나 함께
인증된 비밀 채널로 수신자에게 전달한다.

감사 자료는 공개키 $PK_B$로 암호화하고 비밀키 $SK_B$로 읽는다.
본문의 표기 $pk_{\mathsf{audit}}=PK_B$,
$sk_{\mathsf{audit}}=SK_B$도 같은 키를 뜻한다.
$SK_B$는 Shamir 방식으로 $N$개 비밀 조각에 분산하며,
서로 다른 $K$개 조각으로 복원한다($1\leq K\leq N$).
위원들은 조각을 하나씩 보관한다.~\cite{shamir-secret-sharing-1979}
초기 설정자는 원본 키 사본을 남기지 않고 조각을 각 위원에게
비밀로 전달한다고 가정한다.

외부 승인 뒤 서로 다른 $K$명의 위원이 인증된 비밀 경로로 조각을 전달하면,
감사자는 구성원·배포·키 식별과 공개 확인 자료를 검사하여 $SK_B$를 복원한다.
복원한 키에서 계산한 공개키가 배포된 $PK_B$와 같은지도 확인한다.
감사 승인과 조각 전달은 신뢰된 외부 절차로 수행한다.

유효 키를 복원한 감사자는 원장 기록과 형식을 확인한 뒤
암호문 $ct$를 로컬에서 다음과 같이 복호화한다.
\begin{equation}
 M\leftarrow\mathsf{AuditDec}(sk_{\mathsf{audit}},ct)
 \label{eq:decrypt}
\end{equation}
키 식별·형식·복호화가 맞지 않으면 오류로 처리하고,
$SK_B$와 위원 조각은 공개 거래나 감사 결과에 넣지 않는다.
위원회는 키 복원에 협력하며 기록마다 다시 참여하지 않는다.

배포마다 키 하나를 사용하며, 복원한 키의 열람 범위는 그 키로
암호화한 기록 전체다. 시작점·시점은 업무상 조사 범위를 정한다.
비인가 관측자의 키·조각 보유 조건과 세 보안 성질은
\ref{sec:analysis}절과 부록~\ref{app:security-details}에서 정의한다.

\subsection{개방값·식별자·정수 범위}
\label{app:formats}

해시 입력은 타입과 필드 경계를 보존하는 정규 직렬화를 사용한다.
$H_{\mathsf{tag}}$는 \texttt{zkDPP:tag:v2} 도메인의 해시이며,
\texttt{Note}, \texttt{Voucher}, \texttt{Nullifier} 등 목적을 구분한다.
개방 난수·소유 비밀의 유효 영역, 문자열 정규화와 유한체로의 해시
변환은 배포 설정에서 고정한다.

회계 증표와 이전 증표의 개방값·커미트먼트는 다음과 같다.
\begin{align}
 n&=(d,\rho,q,a,e,o,\omega_N),\nonumber\\
 cm&=H_{\mathsf{Note}}(d,\rho,q,a,e,o,\omega_N)
 \label{eq:note}
\end{align}
\begin{align}
 v&=(d,\rho,q,a,e,o_S,o_R,D,\omega_V),\nonumber\\
 rv&=H_{\mathsf{Voucher}}(d,\rho,q,a,e,o_S,o_R,D,\omega_V).
 \label{eq:voucher}
\end{align}
$d,\rho,q,a,e$는 문서·역할·회계 값이고,
$o,o_S,o_R$는 소유자·송신자·수신자다.
$\omega_N,\omega_V$는 비밀 개방 난수, $D$는 절대 회수 기한이다.
회계 증표의 소비 식별자는 식~\eqref{eq:nf}로,
이전 증표의 공통 해결 식별자는 다음과 같이 계산한다.
\begin{align}
 s_{\mathrm{res}}&=H_{\mathsf{VoucherResolutionSecret}}(\omega_V),\nonumber\\
 rvnf&=H_{\mathsf{VoucherNullifier}}(rv,s_{\mathrm{res}})
 \label{eq:rvnf}
\end{align}
인수와 회수는 각자의 소유권 지식을 증명하지만 동일한 $rvnf$를 쓴다.

문서 해시는 다음 세 필드에 한정한다.
\begin{equation}
 \mathsf{DocumentInfo}=(\mathit{ProductName},\mathit{LotID},\mathit{Unit})
 \label{eq:doc}
\end{equation}
$d=\mathsf{HashDocumentInfo}(\mathsf{DocumentInfo})$이며,
지정 DPP 경로에서 추출한 필드의 누락·잘못된 타입·모호한 인코딩은 거부한다.
객체 조회 키는 \texttt{ObjectRef} 도메인에 타입과 원시 참조를 함께
결합한다. 소비·상태 조회의 $nf,rvnf$도 타입별로 구분한다.
Claim은 등록 여부만 두며 포함 관계 트리·미래 식별자·상태가 없다.

모든 회계 값과 중간 합·곱·나머지는 지정된 정수 범위에 있어야 한다.
유한체에서 값이 돌아오는 계산을 정수 등식으로 오인하지 않도록
범위를 함께 검사한다. 재활용 크레딧의 바닥 함수 배분은
다음 몫·나머지 조건으로 확인한다.
\begin{equation}
 \begin{gathered}
 y=\left\lfloor\frac{xq}{Q}\right\rfloor,\qquad Q>0,\\
 xq=Qy+r,\qquad 0\leq r<Q.
 \end{gathered}
 \label{eq:floor}
\end{equation}
탄소의 비례 배분은 정확히 나누어떨어진다는 가정 아래
$r=0$, 따라서 $xq=Qy$로 표현한다.
분모 0은 허용하지 않는다. 비교식의 단위와 고정 배율도 일치해야 한다.

소비 식별자의 불변성은 효율적인 실행에서 얻은 같은 타입 참조의
개방 증인 $w,w'$에 대해
$\mathsf{SpendId}(w)=\mathsf{SpendId}(w')$이며
생성 시 암호화한 값과도 같다는 조건이다.
이 조건은 개방값의 결합성과 소유자 해시의 충돌 저항성 아래 성립한다.
$\mathsf{OwnerKey}(s)=H_{\mathsf{Owner}}(s)$는 정규 비밀값에
소유자용 해시를 적용한다. 생성·소비 때 소유자 표시가 같은데
서로 다른 비밀을 사용했다면 소유자 해시의 충돌 쌍을 얻는다.
그 사건 밖에서는 두 비밀과 소비 식별자가 같다.
충돌 상계와 전체 논증은 부록~\ref{app:soundness-details}에 둔다.

\subsection{설정과 거래 수락}
\label{app:acceptance}

거래 종류 $k$의 공개 문장은
$x=(u,\boldsymbol\nu,\boldsymbol r,ct)$이고,
비공개 증인은
$w=(\boldsymbol z,\boldsymbol{sk},\boldsymbol{path},\eta)$다.
$u$는 배포·원장·이벤트 경로·정책·권한·루트·기한 등 공개 메타데이터,
$\boldsymbol z$는 실제 입출력 개방값이다.

감사 공개키 $PK_B$는 회로의 고정값이다.
키 식별과 회로·검증기의 대응은 배포 설정에서 정한다.
감사키의 분산·복원은 부록~\ref{app:audit-key}의 절차를 따른다.
이벤트 증인에 $SK_B$나 위원회 비밀 조각을 요구하지 않는다.

부모의 원시 참조 $P_i$와 소비 가능한 출력의 미래
식별자 $\nu_i^{\rm out}$를 한 목록으로 만든다.
$cm,rv$는 정규 유한체 원소이며 타입은 거래 관계로 정한다.
부모가 $m$개이고 소비 가능한 출력이 $s$개이면
\begin{equation}
 M=(P_0,\ldots,P_{m-1},\nu_0^{\rm out},\ldots,\nu_{s-1}^{\rm out}).
 \label{eq:audit-vector}
\end{equation}
이 목록은 실제 입출력 개방값에서 계산한 $M(w)$다.
암호문 $ct=(R,\boldsymbol C)$는 하나의 난수 $0<\eta<\ell$로 만든다.
$\ell$은 생성점 $G$의 위수이며, 아래 덧셈·뺄셈은 평문 유한체에서 한다.
\begin{equation}
 \begin{aligned}
 R&=\eta G,\qquad Z=\eta PK_B,\\
 \kappa&=H_{\mathsf{AuditKey}}(Z_x,Z_y),\\
 C_j&=M_j+H_{\mathsf{AuditMask}}(\kappa,j).
 \end{aligned}
 \label{eq:audit-vector-encryption}
\end{equation}
인덱스 $j$는 전체 목록에서 0부터 시작한다. 저장할 때는 부모의
암호화 값 $m$개와 미래 식별자의 암호화 값 $s$개를 두 배열로
나누지만, 같은 $R$을 쓰는 한 암호문이다. 감사자는 두 배열을 다시
이어 붙이고 $Z=SK_B R$로 같은 $\kappa$를 계산해 각 마스크를 뺀다.
미래 식별자 배열에서도 인덱스를 0으로 다시 시작하지 않는다.

Transfer의 출력 순서는 이전 증표·잔여 회계 증표이고,
평문 목록은 $(cm_{\rm in},rvnf_T,nf_C)$다.
Claim에는 미래 식별자가 없고, Exit·Issue는 부모만 암호화한다.
Entry는 부모 없이 출력 Note의 미래 식별자만 암호화한다.
따라서 아홉 거래 모두 암호화할 값이 적어도 하나 있다.

키 유도 입력은 식~\eqref{eq:audit-vector-encryption}의 점 좌표이며,
거래 문맥·목록 길이·타입은 회로의 정해진 형상과 전체 공개 문장으로 결합한다.
계약은 공개 입출력 뒤에 $R$의 좌표·두 암호화 배열을 같은 순서로
붙여 증명을 검증하고, 검증한 자료를 그대로 기록해야 한다.
회로는 난수의 범위와 비영 값을 검사하며, 거래 사이의 독립 난수
생성은 클라이언트의 책임이다. 동일 키에서 난수를 재사용하면
같은 위치의 마스크가 반복되어 평문 간 차이가 드러날 수 있다.
복호화 정확성은 이 관계의 모든 허용 난수에서 $M(w)$가 복원된다는
조건이다.

공통 암호·생애주기 관계와 선택한 회계 관계를 합쳐 검사한다.
\begin{equation}
 \mathcal{R}_k=\mathcal{R}^{\mathrm{core}}_k
                  \land\mathcal{R}^{\mathrm{acct}}_k
 \label{eq:relation}
\end{equation}
공통 관계에는 개방·소유권·포함 관계·소비 식별자·정수 범위,
중복 금지와 출력·감사 결합을 포함한다.
고정 경로는 이벤트별 검증기 선택으로 종류를 결합하고,
등록 경로는 정책 $p$와 필요한 적용 범위를 증명에 결합한다.

배포자는 암호·직렬화·거래 형상·역할·위원회와 고정 검증기를 설치한다.
회계·이전 증표는 별도의 깊이 32 추가 전용 머클 트리에 저장하며,
초기 트리는 비어 있고 그 루트는 수락된 루트다.
노드·리프 수·현재 및 과거 수락 루트를 보관해 포함 경로를
공개 자료로 재구성할 수 있게 한다. 용량 초과는 거래 전체를 거부한다.
소비·생성 색인, Claim과 감사 기록도 빈 상태에서 시작하고
$\aid$는 1부터 증가한다.

등록 정책의 식별자는 거래 종류 $k$, 등록 주체 $a$, 그 주체의
정책 번호 $b$와 버전 $v$로 계산한다.
\begin{equation}
 p=H_{\mathsf{PolicyRef}}(k,a,b,v).
 \label{eq:policy-ref}
\end{equation}
원장은 권한 있는 등록 주체에 $a$를 부여하고 정책 번호·버전을
예약한 뒤, 그 식별자에 검증기와 정확한 입출력 수를 한 번 등록한다.
입출력 수는 불변 등록 정보와 회로에서 고정하며,
거래 때 두 형식이 일치하는지 검사한다.
관계 명세와 회로·검증키의 대응을 공개하고,
가공의 각 $(m,n)$ 형태는 별도 정책 버전 또는 번호와 검증키를 쓴다.
계산 규칙·형상·검증키를 바꿀 때는 새 식별자를 등록한다.
검증기 메타데이터는 불변이고 정책 활성값은 참에서 거짓으로만 바뀐다.
비활성화는 이후 제출에 적용하며 과거 수락·Claim을 취소하지 않는다.
가공의 불투명 식별자 \texttt{scope}에 대해 제출자 또는 자격증명의
범위 권한과 $(scope,p)$ 사용 권한을 함께 검사한다.
소유권 지식만으로 모든 정책을 사용할 수는 없다.
고정 이벤트의 정책 표식 0은 등록 정책 식별자와 구분한다.

증명자는 실제 개방값·포함 경로를 준비하고,
입력 소비 식별자, 출력 참조·미래 식별자와 감사 암호문을 계산하여
전체 관계의 증명을 만든다. 원장은 경로·형상·정책·권한,
수락 루트·현재 미소비·활성·실행 시 기한과 모든 출력 키의 신규성을 검사한다.
같은 거래 안의 실제 입력 참조, 타입별 소비 식별자와 출력 조회 키도
각각 중복될 수 없다. 이미 등록된 Claim의 재생성도 거부한다.
통과하면 한 트랜잭션에서 $\aid$ 배정, 소비 기록,
트리 추가 또는 Claim 등록, 감사 기록과 모든 출력의 생성 색인을
갱신하고 다음 $\aid$로 진행한다. 실패하면 모두 취소한다.
수락된 소비·생성·감사 기록은 불변이다.

상태는 별도 트리가 아닌 타입별 미래 소비 식별자의 매핑으로 관리한다.
기본값은 Active, Revoked는 종단이며 Active에서 바로 갈 수 없다.
감사자는 실제 생성 기록을 확인하여 해당 객체의 미래 소비 식별자를 선택한다.

Entry는 유효한 명시 초기 탄소 $e_0$를 보존한다.
생략만 클라이언트에서 0으로 정규화하며 잘못된 명시 값을 0으로 바꾸지 않는다.
그 값은 개방값·증인에 결합되고, 기본값을 회로의 $e_0=0$ 제약으로
해석하지 않는다. Transfer의 $D$는 증표 안의 값과 같고,
생성 실행 시 현재보다 미래이면서 배포의 허용 범위 안에 있어야 한다.
Recall은 실행 에포크가 $D$인 경우도 허용하고 이후 제출은 거부한다.

\subsection{추적 규칙}
\label{app:trace-details}

실제 그래프는 수락 때의 실제 생성·소비 값으로 정의한다.
객체에서 그것을 소비한 이벤트로, 이벤트에서 전체 출력으로 간선을 잇는다.
생성이 소비보다 앞서므로 순환이 없다.
인증 조회는 고정 $H_t$의 실제 값 또는 오류를 반환하며,
확정되지 않은 블록 번호만으로 확정성을 판단하지 않는다.

선행 객체 시작은 그 객체·생성 이벤트·둘 사이 간선을 포함한다.
선행 기록 시작은 지정 이벤트와 모든 직접 출력·출력 간선을 포함한 뒤
실제 부모를 거슬러 올라간다. 따라서 출력 없는 Exit도 시작점이 된다.
재귀 중에는 부모의 생성 기록을 읽되 그 기록의 다른 출력은 추가하지 않는다.
각 부모의 타입·수·순서·중복과 더 이른 생성 순서를 확인한다.
부모가 비어 있으면 Entry여야 하며 초기 등록 집합에 기록한다.

후속 객체 시작은 그 객체 하나, 기록 시작은 그 기록의 전체 출력을
초기 큐에 넣는다. 시작 기록 자체와 출력 간선을 출발점이라는 이유로
추가하지 않으며, 출력 없는 기록의 후속 결과는 빈 결과다.
이는 종료 이전의 선행 조사와 다르다.
미소비 Note/Voucher는 타입 참조·미래 식별자와 함께 각각 $F_N,F_V$에,
도달 Claim과 Exit는 별도의 종단 집합에 기록한다.

객체의 존재·생성 기록·출력 위치·이벤트 형상과 공개 트랜잭션 자료를
확인하고 전체 벡터를 복호화한 뒤 정해진 위치에서 필요한 값을 얻는다.
소비 식별자 조회가 양수이면 해당 기록이 실제 생성보다 나중이고,
공개 입력에 같은 타입의 식별자가 있는지 확인한다.
소비 이벤트의 부모 목록에 현재 객체가 포함됐는지도 검사한다.
함께 복호화한 다른 부모를 후속 그래프의 선행 경로로 추가하지 않는다.
부모와 소비 식별자의 실제 입력 결합은 수락 증명에 의존한다.

같은 이벤트의 복호화·출력 확장은 한 번만 수행할 수 있지만,
재방문에서 새로 도달한 객체와 입력·출력 간선은 보존한다.
같은 생성 기록을 읽는 것과 그 기록의 모든 출력을 후속으로
추적하는 것은 다른 처리다.
동일 객체의 상충하는 식별자, 서로 다른 객체의 같은 타입 식별자는
중복으로 합치지 않고 오류로 처리한다.
숫자가 같아도 타입이 다른 객체·식별자는 구분한다.

키·자료 부재나 발견된 디코딩·복호화·일관성 오류는 실패다.
부분 진단 자료를 반환할 수 있어도 완성된 그래프나 전체 대상으로
보고할 수 없다. 승인 부재·잘못된 원장 문맥은 집행 전에 거부한다.

\subsection{감사·차단의 집행과 반환}
\label{app:freeze-execution}

\ref{sec:audit}절의 절차는 추적·일관성 검사를 마친 뒤 동결을 시작한다.
검사가 실패하면 \Alg{TRACE\_FAILED}로 종료하고 이 실행의 동결 거래는
제출하지 않는다. 추적이 성공하면 전체 미소비 집합을 $T$로 고정하며,
비어 있으면 \Alg{NO\_LIVE\_TARGETS}를 반환한다.

대상별 처리는 Table~\ref{tab:freeze}를 따른다.
한 대상의 실패가 이미 성공한 동결을 되돌리지는 않는다.
가능한 다른 대상은 계속 처리하되 시도·확인은 유한 한도 안에서 끝낸다.
미시도·실패·불명확한 제출은 사유와 함께 남긴다.
\begin{table}[htbp]
\centering\small
\caption{Per-target action at execution time}
\label{tab:freeze}
\begin{tabularx}{\columnwidth}{@{}YY@{}}
\toprule
Observed state & Action in this run \\
\midrule
Unspent, Active & Submit a typed Freeze transaction \\
Unspent, Frozen/Revoked & Record an existing block \\
Already spent & Retain target and consumer as a race \\
Read failure & Retain target as unconfirmed \\
\bottomrule
\end{tabularx}
\end{table}

거래 영수증은 제출 결과의 근거이고 최종 차단 여부는 공통 $H_r$에서 판단한다.
동결 제출이 다른 실행의 선행 동결 때문에 거부됐더라도,
공통 시점에서 이미 차단됐음이 확인될 수 있다.
먼저 소비된 사실을 관찰했다면 그 소비 이전의 시점을 골라 완료로
보고할 수 없다. 결과에는 시작점, $H_t$, 추적 결과, 전체 $T$,
대상별 시도·영수증·상태, 공통 $H_r$ 또는 그 부재와 최종 판정을 남긴다.
정상 추적과 비어 있지 않은 $T$를 얻은 실행은 \ref{sec:audit}절의
세 완료 조건을 모두 만족하면 \Alg{COMPLETE\_AT\_CHECKPOINT}를,
그 밖에는 전체 대상을 보존한 \Alg{INCOMPLETE}를 반환한다.
